\documentclass[11pt]{article}
\usepackage[worksheet]{la}
\lamodule{1}{5}{Solution Sets of Linear Systems}

\begin{document}

\section{Solution Sets of Linear Systems}

{\Large\bfseries Activity Worksheet: Describing All the Solutions}\par

\textbf{Names:} \rule{2.4in}{0.4pt}\quad\rule{2.4in}{0.4pt}\\[6pt]
\hspace*{3.8em}\rule{2.4in}{0.4pt}\quad\rule{1.3in}{0.4pt}\quad\textbf{Date:} \rule{0.9in}{0.4pt}

\begin{roles}
Work in groups of three or four. Rotate the roles after each activity.
\begin{itemize}[nosep]
  \item \textbf{Facilitator:} keeps the group on task and makes sure everyone can explain the answer.
  \item \textbf{Recorder:} writes the group's work on \emph{this} sheet.
  \item \textbf{Checker:} verifies every claimed solution by multiplying it by the matrix.
  \item \textbf{Reporter:} shares the group's reasoning with the class.
\end{itemize}
Everyone must be able to explain \emph{why} an answer is correct, not just what it is. If you finish
early, move on to the extension activities.
\end{roles}

% =====================================================================
\begin{activity}{Activity 1: The homogeneous equation}{about 14 minutes}
Let
\[
A=\begin{bmatrix}1&2&3&-1\\2&5&8&-4\\-1&-1&-1&-1\end{bmatrix}.
\]
\begin{enumerate}[label=(\alph*),itemsep=6pt]
  \item Row reduce $[A\ \bv{0}]$ to reduced echelon form. Circle the pivot positions. Which variables are free?
        \workspace{1.4in}
  \item Write the solution set of $A\bv{x}=\bv{0}$ in parametric vector form $\bv{x}=x_3\bv{u}+x_4\bv{v}$. Write down $\bv{u}$ and $\bv{v}$.
        \answerline
  \item \textbf{Checker:} verify $A\bv{u}=\bv{0}$ and $A\bv{v}=\bv{0}$ by computing each product.
        \answerline
  \item Let $\bv{w}=2\bv{u}-3\bv{v}$. Compute $\bv{w}$. Without computing $A\bv{w}$ directly, explain why $A\bv{w}=\bv{0}$.
        \answerline
  \item Is $(1,1,1,1)$ a solution? Compute $A(1,1,1,1)$ to find out.
        \answerline
\end{enumerate}
\textbf{Think about it.} A $3\times4$ homogeneous system can have how many free variables: what is the smallest possible number, and what is the largest? Give a matrix for each extreme.
\answerline\answerline
\end{activity}

\newpage
\begin{activity}{Activity 2: Shifting the solution set}{about 14 minutes}
Use the same matrix $A$ as in Activity~1 and let $\bv{b}=(3,7,-2)$.
\begin{enumerate}[label=(\alph*),itemsep=6pt]
  \item Row reduce $[A\ \bv{b}]$. Is the system consistent?
        \workspace{1.4in}
  \item Write the solution set of $A\bv{x}=\bv{b}$ in parametric vector form $\bv{x}=\bv{p}+x_3\bv{u}+x_4\bv{v}$. Compare the vectors $\bv{u},\bv{v}$ with Activity~1.
        \answerline
  \item Verify that your $\bv{p}$ is a solution by writing $A\bv{p}$ as a combination of columns of $A$. Which columns?
        \answerline
  \item Choose values for $x_3$ and $x_4$ to produce a second solution $\bv{q}$. Compute $\bv{q}-\bv{p}$ and check that it solves $A\bv{x}=\bv{0}$.
        \answerline
  \item Is $\bv{0}$ a solution of $A\bv{x}=\bv{b}$? How does this show that the solution set is \emph{not} a span?
        \answerline
  \item \textbf{Warm-down in $\R^2$.} The equation $x_1-2x_2=3$ has solution set $\bv{p}+x_2\bv{d}$. Find $\bv{p}$ and $\bv{d}$, and sketch the points $\bv{p}$ and $\bv{p}+\bv{d}$ in the space below. Where is the homogeneous solution set $\Span\{\bv{d}\}$ in relation to it?
        \workspace{1.6in}
\end{enumerate}
\textbf{Think about it.} Suppose a different right-hand side $\bv{b}'$ also gives a consistent system with the same $A$. What would be the same and what would be different about its solution set?
\answerline\answerline
\end{activity}

\textbf{Exit ticket} (individually, before leaving). In one or two sentences: if you are handed one solution of $A\bv{x}=\bv{b}$, what else do you need to know to describe \emph{all} solutions? What question do you still have?
\answerline\answerline

\newpage
% =====================================================================
\subsection*{Extension activities \quad (if you finish early, or as take-home)}

\begin{ext}{Activity 3: A three-junction road network}{}
Three junctions $A,B,C$ lie on a one-way triangle $A\to B\to C\to A$. Let $x_1,x_2,x_3$ be the traffic (cars per hour) on $A\to B$, $B\to C$, $C\to A$. Cars enter at $A,B,C$ at the rates $40,20,30$ and leave at the rates $10,50,30$.
\begin{enumerate}[label=(\alph*),itemsep=6pt]
  \item Write the balance equation (flow in $=$ flow out) at each junction and put the system in the form $M\bv{x}=\bv{b}$.
  \item Check that the total entering equals the total leaving. Why is this necessary?
  \item Solve the system in parametric vector form. What is the meaning of the direction vector?
  \item Since traffic counts cannot be negative, which values of the free variable are allowed?
  \item A traffic counter on $B\to C$ reads $12$ cars per hour. Find the other two flows.
\end{enumerate}
\workspace{1.8in}
\end{ext}

\begin{ext}{Activity 4: A plane in $\R^3$}{}
Let $\bv{p}=(1,0,2)$, $\bv{s}=(1,-1,0)$, $\bv{t}=(0,1,1)$. The points $\bv{p}+a\bv{s}+b\bv{t}$ ($a,b\in\R$) form a plane.
\begin{enumerate}[label=(\alph*),itemsep=6pt]
  \item Find one linear equation $c_1x_1+c_2x_2+c_3x_3=d$ satisfied by $\bv{p}$, $\bv{p}+\bv{s}$ and $\bv{p}+\bv{t}$.
  \item Show that $\bv{s}$ and $\bv{t}$ are solutions of the corresponding homogeneous equation, and that neither is a multiple of the other.
  \item Explain why the solution set of your equation is exactly the plane (count free variables).
\end{enumerate}
\workspace{1.8in}
\end{ext}

\begin{ext}{Activity 5: Two solutions are enough}{}
Suppose $A\bv{x}=\bv{b}$ has two different solutions $\bv{p}$ and $\bv{q}$. Show that it has infinitely many, by exhibiting a one-parameter family of solutions built from $\bv{p}$ and $\bv{q}$. For which values of $c$ is $c\bv{p}+(1-c)\bv{q}$ a solution? Is $\bv{p}+\bv{q}$ a solution when $\bv{b}\neq\bv{0}$?
\workspace{1.6in}
\end{ext}

\end{document}
